\documentclass[journal]{IEEEtran}
\usepackage{amsmath,amssymb,booktabs,url}

\title{EvoGraph: Evidence-Aware Static--Dynamic Malware Representation}
\author{EvoGraph Project}

\begin{document}
\maketitle

\begin{abstract}
EvoGraph specifies an evidence-aware integration of portable-executable
features and sandbox telemetry for binary malware classification. Its
methodological contribution is a six-type schema combining source-order
behavioral windows, typed resource interactions, explicit missingness, and
a controlled flat-vector comparison. Standard relational message passing
and type-specific attention integrate these representations; no new graph
operator is claimed. This article describes implemented methodology and
an evaluation protocol, not empirical detection results.
\end{abstract}

\begin{IEEEkeywords}
Malware representation, heterogeneous graphs, multimodal learning.
\end{IEEEkeywords}

\section{Methodological Scope}
Static features summarize executable structure, whereas sandbox reports
describe activity observed under a particular execution environment.
Flattening these observations removes process ownership and local order;
constructing overly connected graphs can instead invent relationships.
EvoGraph integrates both modalities while distinguishing observed
execution, inventory evidence, and unavailable telemetry. The contribution
is this evidence-aware representation and its comparison design, rather
than an established performance improvement or a claim of first use.

Inputs are user-supplied CAPE JSON reports~\cite{cape}, binary labels, and
optional static sidecars or binaries for read-only parsing. Sandbox
execution and CAPE API transport are outside scope; no API client or
sandbox is claimed here.

\section{Evidence and Flat Representation}
Preparation preserves normalized evidence, source pointers, report hashes,
metadata, and complete original EMBER dictionaries; original reports remain
external and must be retained. Available executable hashes are checked
across inputs, conflicting identities or labels are rejected, and duplicate
identities are excluded. Metadata alone cannot verify report--binary
correspondence when the report omits its target hash. Labels, sandbox
verdicts, and signatures are not model features.

The static vector has 2,381 dimensions from official EMBER feature version
2~\cite{emberrepo}. Anderson and Roth's original paper describes EMBER2017
feature version 1~\cite{ember}, not this later schema; its 2,351-dimensional
vectors are not padded. The flat comparator concatenates static features
with a 2,048-dimensional dynamic representation (Table~\ref{tab:blocks}),
giving 4,429 dimensions. D1 is deterministic; D2 uses stable signed BLAKE2b
hashing of normalized executed-API frequencies. D3--D10 are learned jointly
with the classifier, making exported dynamic vectors checkpoint-dependent,
not pretrained features.

\begin{table}[t]
\caption{Flat Dynamic Representation: 2,048 Dimensions}
\label{tab:blocks}
\centering
\small
\setlength{\tabcolsep}{4pt}
\begin{tabular}{@{}llr@{}}
\toprule
Block & Encoding / aggregation & Dim.\\
\midrule
D1 Statistics & Observations and availability & 128\\
D2 API identity & Signed frequency hash & 512\\
D3 API sequences & BiGRU; event/stream attention & 256\\
D4 Files & Path CharCNN; mean/max & 256\\
D5 Processes & Name/command; mean/max & 256\\
D6 Registry & Hierarchical BiLSTM; mean/max & 128\\
D7 Network & Identifier CharCNN; mean/max & 128\\
D8 Other resources & Typed CharCNN; mean/max & 128\\
D9 Arguments & Typed semantics; mean/max & 128\\
D10 Payloads & Mask-aware static; mean/max & 128\\
\bottomrule
\end{tabular}
\end{table}

Unavailable sections differ from observed empty sections. Missing static
vectors remain null in prepared data; masks and learned availability
embeddings accompany numerical placeholders, rather than presenting zeros
as measured features. D1 retains coverage flags. Independently, the graph's
binary node receives 14 flags: 12 section-coverage indicators and separate
trace/window truncation indicators. Payloads retain per-vector static
availability and report-level coverage.

\section{Evidence-Aware Heterogeneous Graph}
Each report yields exactly six node types: \emph{binary, process, API,
window, resource, payload}. The binary node encodes static features through
$2381\!\to\!512\!\to\!256$, adding availability and projected telemetry.
Process inputs concatenate a 64-dimensional name CharCNN, a 32-dimensional
command projection, and 32 local statistics. API inputs combine identity
(64), category (8), and imported/resolved flags (2), projected to 128.
Windows use a BiGRU with 128 units per direction,
event attention at 256 dimensions, and projection to 128. Resource nodes
combine kind-specific string encoders, type embeddings, and local evidence
statistics at 128 dimensions. Payloads use mask-aware
$2381\!\to\!64\!\to\!128$ projections. All types then project to hidden
width 256.

Windows contain at most 32 retained source-call rows within one
process/thread; unidentified threads form a separate stream per process.
Source order is preserved without global timestamp sorting. Repeated calls
remain compressed rows. Imports and resolved APIs identify capabilities,
not executions; count-only \texttt{apistats} can supply frequencies but
never fabricate call sequences.

Resource identity is keyed separately by kind and observed string.
Encoder-only masking of local usernames and literal IP endpoints does not
merge graph identities. Resource edges require explicit names, paths, or
pretty values: numeric API handles are never guessed into bindings.
Summary/network inventories do not invent process ownership or event edges.

The 17 forward relation triples are binary-to-process \emph{launches},
binary-to-API \emph{imports}, binary-to-payload \emph{produces};
process-to-process \emph{spawns}, process-to-API \emph{calls},
process-to-window \emph{executes}, process-to-resource \emph{loads/accesses},
process-to-payload \emph{drops}; window-to-API \emph{contains\_api},
window-to-window \emph{next\_window}, and window-to-resource
\emph{reads, writes, deletes, modifies, connects\_to, accesses}.
Each has an exact reverse with independent parameters, yielding 34 directed
relation types.

Static imports and runtime windows meet at shared API-symbol nodes, exposing
declared-versus-observed behavior without equating imports with calls.
Separating symbols from sequence windows allows the model to combine API
identity, local order, and resource context rather than duplicating one
whole-report embedding. This cross-view, evidence-constrained construction
is the central methodological hypothesis to be tested.

\section{Relational Integration}
Three relation-mean R-GCN-style layers~\cite{rgcn} operate at width 256.
For incoming neighbors $\mathcal N_r(v)$, node type $t=t(v)$, and layer
$\ell$, the implemented update is
\begin{align}
m_{v,r}^{(\ell)}
&=\frac{\sum_{u\in\mathcal N_r(v)}W_r^{(\ell)}h_u^{(\ell)}}
{\max(1,|\mathcal N_r(v)|)},\label{eq:mean}\\
z_v^{(\ell)}
&=W_{0,t}^{(\ell)}h_v^{(\ell)}+b_{0,t}^{(\ell)}
+\sum_r m_{v,r}^{(\ell)},\notag\\
h_v^{(\ell+1)}
&=\operatorname{LN}_{t}^{(\ell)}\!\left(h_v^{(\ell)}+
\operatorname{Drop}_{0.2}\!\left(\operatorname{GELU}(z_v^{(\ell)})\right)\right).
\label{eq:residual}
\end{align}
Empty sums contribute zero. Means are normalized per relation, then summed,
not divided by total degree or the number of relations. PyG implements
this with relation-specific mean SAGE aggregators without root transforms,
bias, or normalization, summed by \texttt{HeteroConv}; self transforms,
residuals, and layer normalization are type-specific. This instantiates
established relational aggregation, not a new GNN operator.

Let $H_v=h_v^{(3)}$ and $V_t$ denote nodes of non-binary type $t$.
Type-specific attention and fusion are
\begin{align}
\alpha_{v,t}&=\frac{\exp(a_t^\top H_v)}
{\sum_{u\in V_t}\exp(a_t^\top H_u)},\notag\\
p_t&=\mathbf{1}_{|V_t|>0}\,P_t\!\left(\sum_{v\in V_t}
\alpha_{v,t}H_v\right)+c_t,\label{eq:pool}\\
g&=\phi_{512}\!\left([H_b\Vert p_{\rm proc}\Vert p_{\rm api}
\Vert p_{\rm win}\Vert p_{\rm res}\Vert p_{\rm pay}]\right),\notag\\
\widehat p&=\sigma\!\left(f_{576\to256\to64\to1}
([g\Vert\psi_{128\to64}(s)])\right).\label{eq:fusion}
\end{align}
Attention is evaluated only for nonempty types. Here $P_t$ is an affine
256-to-128 projection; empty pools are zero before
$c_t$, which is a learned coverage embedding only for payloads and zero
otherwise. Thus $g$ projects $256+5(128)=896$ dimensions to 512.
The maps $\phi$ and $\psi$ use affine--GELU--dropout and affine--GELU,
respectively; $f$ uses GELU/dropout between linear layers.
The hybrid appends only projected D1 statistics $s$, producing 576
dimensions: it neither computes nor appends the full 2,048-dimensional
dynamic branch. Graph-only removes D1, retaining missingness and local
statistics; static-only and dynamic-only are additional ablations.

\section{Protocol and Current Status}
Default stratified splits share approximately 60/20/20 train/validation/test
partitions across modes, with both labels required in every partition.
Non-mask D1 values receive $\log(1+x)$ and training-only standardization;
static moments use present training sample vectors, also scaling payloads.
Masks remain unchanged. Weighted binary cross-entropy and AdamW update
encoders only on training samples. Validation loss controls early stopping;
validation F1 selects the decision threshold. Test data neither fit
preprocessing nor select checkpoints or thresholds. The comparison controls
partitions and provenance, not parameter count or receptive field.

Family-disjoint evaluation requires explicit groups for every record,
including benign application annotations. Chronological evaluation requires
genuine observation datetimes and rejects tied boundaries or one-class
partitions; CAPE execution time is not sample first-seen time. Extraction
provenance must record the official EMBER revision and dependency profile,
used consistently for training and inference; the local LIEF~0.12.3 profile
is not an exact EMBER2018/LIEF~0.9.0 replication.

No trained research model or full multimodal benchmark is currently
available. Existing tests and smoke checks exercise plumbing only, not
detection effectiveness. Defaults cap retention at 20,000 source rows and
256 windows;
window retention favors process/thread prefixes, inventories can exceed
retained traces, and caps do not bound graph memory. Label quality,
missingness, and sandbox/source confounding require auditing before empirical
claims. HGT comparisons and incremental learning are deferred.

\begin{thebibliography}{4}
\bibitem{cape}
CAPE Sandbox, ``CAPE Sandbox Book,'' documentation. [Online]. Available:
\url{https://capev2.readthedocs.io/en/latest/}

\bibitem{emberrepo}
Elastic, ``Elastic malware benchmark for empowering researchers,'' official
repository, revision \path{d97a0b523de02f3fe5ea6089d080abacab6ee931}.
[Online]. Available: \url{https://github.com/elastic/ember}

\bibitem{ember}
H.~S. Anderson and P.~Roth, ``EMBER: An open dataset for training static PE
malware machine learning models,'' arXiv:1804.04637, 2018.

\bibitem{rgcn}
M.~Schlichtkrull, T.~N. Kipf, P.~Bloem, R.~van den Berg, I.~Titov, and
M.~Welling, ``Modeling relational data with graph convolutional networks,''
arXiv:1703.06103, 2017.
\end{thebibliography}

\end{document}
